\section{Conclusão da Obra: O Legado Preditivo}

Este livro percorreu uma rigorosa trajetória ontológica e operacional, mapeando a metamorfose da odontologia analógica para a vanguarda dos Gêmeos Digitais. A tese irrefutável que emerge de nossas investigações é que o Gêmeo Digital não é uma ferramenta satélite, mas o novo paradigma gravitacional em torno do qual toda a prática, o ensino e a pesquisa odontológica orbitarão nas próximas décadas.

A odontologia repousava na habilidade de tratar a patologia instalada. Com o Gêmeo Digital, ela ascende à capacidade de prever a falha antes que ocorra no plano biológico. Contudo, essa convergência entre \textit{Machine Learning}, modelagem in silico e biomecânica só alcançará sua plenitude ética se ancorada em ecossistemas colaborativos e de código aberto, como o \textbf{OpenCode Ecosystem}, que impedem a monopolização do conhecimento \cite{robles2023opentwins}. 

O Brasil, munido da força capilar do SUS e de sua histórica capacidade de inovação na saúde pública e bioinformática, possui o lastro estrutural necessário para liderar essa transição. A odontologia do futuro será inevitavelmente digital, inegavelmente preditiva e, se agirmos com sabedoria legislativa e ética (\autoref{ch:etica}), universalmente participativa. Os algoritmos estão treinados e as malhas tridimensionais, construídas. O desafio agora é, sobretudo, de governança e coragem clínica.
